% Generated draft: requires mathematical and source review.
\documentclass[11pt]{article}
\usepackage[margin=25mm]{geometry}
\usepackage{fontspec}
\setmainfont{Latin Modern Roman}
\usepackage{amsmath,amssymb,mathtools,unicode-math}
\setmathfont{Latin Modern Math}
\usepackage{enumitem,needspace,xurl,hyperref}
\setlength{\parindent}{0pt}
\setlength{\parskip}{5pt}
\setlength{\emergencystretch}{3em}
\newcommand{\N}{\mathbb{N}}
\newcommand{\Z}{\mathbb{Z}}
\newcommand{\Q}{\mathbb{Q}}
\newcommand{\R}{\mathbb{R}}
\newcommand{\C}{\mathbb{C}}
\newcommand{\F}{\mathbb{F}}
\newcommand{\A}{\mathbb{A}}
\newcommand{\PP}{\mathbb{P}}
\newcommand{\E}{\mathbb{E}}
\newcommand{\eps}{\varepsilon}
\newcommand{\dd}{\,\mathrm{d}}
\begin{document}
\section*{UnsolvedMath ID: 30005902}
\textbf{Source problem code:} catalogue entry
\subsection{Definitions and mathematical statement}
The supplied record does not specify the cohomology theory, coefficients, ground field, or construction of the graded Lie bracket. The following is therefore a conditional formulation under the usual Hopf-cohomology interpretation; these choices require confirmation from the cited report.

Let \(k\) denote the unspecified ground field, and let \(H\) be a Hopf algebra over \(k\), with coproduct \(\Delta\) and counit \(\varepsilon\). Give \(k\) the trivial left \(H\)-module structure \(h\cdot a=\varepsilon(h)a\), and interpret cohomology as
\[
\mathrm{H}^{n}(H,k)=\operatorname{Ext}^{n}_{H}(k,k),
\qquad
\mathrm{H}^{>0}(H,k)=\bigoplus_{n\geq 1}\mathrm{H}^{n}(H,k).
\]
Interpret the bracket as the canonical Gerstenhaber bracket on this cohomology, with degree convention
\[
[-,-]\colon \mathrm{H}^{p}(H,k)\times\mathrm{H}^{q}(H,k)
\longrightarrow \mathrm{H}^{p+q-1}(H,k).
\]
Thus the bracket has degree \(-1\), and becomes a degree-zero graded Lie bracket when a class of cohomological degree \(n\) is assigned shifted degree \(n-1\).

Let \(H\text{-}\mathrm{Mod}\) denote the category of left \(H\)-modules, with tensor product over \(k\), unit \(k\), and action
\[
h\cdot(m\otimes n)=\sum h_{(1)}m\otimes h_{(2)}n,
\qquad \Delta(h)=\sum h_{(1)}\otimes h_{(2)}.
\]
A braiding means a natural family of invertible \(H\)-linear maps \(c_{M,N}\colon M\otimes_k N\to N\otimes_k M\) satisfying the two hexagon identities for this monoidal structure.

Under these interpretations, the question is whether every such \(H\) for which \(H\text{-}\mathrm{Mod}\) admits no braiding satisfies
\[
\forall p,q\geq 1\quad
\forall \alpha\in\mathrm{H}^{p}(H,k)\quad
\forall \beta\in\mathrm{H}^{q}(H,k),
\qquad [\alpha,\beta]=0.
\]
The range of permitted ground fields and any additional hypotheses on \(H\) remain unspecified. The record's title and literature assessment use ``non-quasitriangular,'' whereas its question uses absence of a braiding on the module category. Their intended relationship, including the precise module category considered, requires source confirmation. If the intended cohomology is Hochschild cohomology or uses other coefficients, the displayed formulation must be revised.
\subsection{Short English statement}
Under the provisional interpretation \(\mathrm{H}^{*}(H,k)=\operatorname{Ext}^{*}_{H}(k,k)\), ask whether the Gerstenhaber bracket of every pair of positive-degree cohomology classes vanishes whenever the usual tensor category of \(H\)-modules admits no braiding. The cohomology, bracket, ground-field restrictions, and intended non-quasitriangular hypothesis require confirmation. The supplied record labels the problem open; this draft does not independently verify that status.
\subsection{Sources}
\begin{itemize}
\item Local UnsolvedMath problem registry, permanent ID \(30005902\), exact source problem code \texttt{catalogue entry}, title \emph{Vanishing Lie Brackets for Non-Quasitriangular Hopf Algebras}. The record assigns the problem to the 2024 Oberwolfach Reports collection and names no proposer.
\item Supplied source citation: \emph{Hochschild (Co)Homology and Applications} (2024), Oberwolfach Reports, DOI 10.4171/owr/2024/20, pp.\ 1097--1162. \url{https://doi.org/10.4171/owr/2024/20}. The record says the official report was checked and the question isolated with repaired notation; this is supplied verification metadata, not an independent verification in this draft. The supplied original-statement field contains only introductory context, without an explicit mathematical question.
\item Supplied literature source: Hermann, monoidal categories and the Gerstenhaber bracket. \url{https://doi.org/10.1090/memo/1151}.
\item Supplied literature source: Karadag, Gerstenhaber brackets on Hopf algebras. \url{https://doi.org/10.1016/j.jpaa.2021.106903}.
\item Supplied literature source: Oberwolfach Report 20/2024. \url{https://doi.org/10.4171/owr/2024/20}.
\item Literature assessment dated exactly \texttt{2026-08-22}: the supplied curation labels the problem open and states that the bracket vanishes in braided cases and several non-quasitriangular examples, while universal vanishing for non-quasitriangular Hopf algebras remains unknown. This is supplied dated curation derived from the OWR record, not independently verified literature assessment; its status and cited claims require independent checking.
\end{itemize}
\end{document}
